\section{Arm CCA / RME / RMM 基础架构}

\begin{frame}[t,shrink=6]{Arm CCA / RME / RMM 基础架构：方向开场}
\DeckKicker{方向开场}
\SlideMeta{证据等级}{方向综述}{来源状态：公开材料综合}
{\large\bfseries\color{Ink} Arm CCA / RME / RMM 基础架构 的核心是先界定保护对象、管理边界和证据强度，再用三篇主讲材料串起技术演进。\par}\vspace{1.8mm}
\begin{columns}[T,totalwidth=\textwidth]
  \begin{column}{0.45\textwidth}
    \fcolorbox{Rule}{Paper}{%
  \begin{minipage}[t]{0.97\linewidth}
    \vspace{1.1mm}
    \hspace{1.4mm}\begin{minipage}{0.92\linewidth}
      \PanelTitle{内容说明}
      \scriptsize \begin{itemize}
  \item 围绕 Realm、RME、RMM、granule ownership 与 lifecycle 建立 Arm CCA 主线。
  \item 固定三篇主讲材料；`wu2024rmm` 作为 RMM verification auxiliary，不计入三篇主讲。
  \item 先讲方向痛点和设计轴，再逐篇说明贡献、限制、证据强度和可落地条件。
  \item 对规范、Survey、draft、vendor evidence 保持显式标注，避免把背景材料写成一手系统证明。
\end{itemize}
    \end{minipage}
    \vspace{1mm}
  \end{minipage}}
  \end{column}
  \begin{column}{0.49\textwidth}
    \fcolorbox{Rule}{Panel}{%
  \begin{minipage}[t]{0.97\linewidth}
    \vspace{1.1mm}
    \hspace{1.4mm}\begin{minipage}{0.92\linewidth}
      \PanelTitle{三篇主讲材料的分工}
      \scriptsize {\tiny
\begin{tabularx}{\linewidth}{@{}YYY@{}}
\toprule
\textbf{论文} & \textbf{定位} & \textbf{证据边界} \\
\midrule
Arm Confidential Computing Architecture & 开创/基础入口 & E1 / 基础证据 \\
Arm CCA 架构规范 & 代表性改进 & E0 / Spec-standard SOTA \\
Realm Management Monitor 规范 & 当前 SOTA/补充边界 & E0 / Spec-standard SOTA \\
\bottomrule
\end{tabularx}
}
    \end{minipage}
    \vspace{1mm}
  \end{minipage}}
  \end{column}
\end{columns}
\vspace{1mm}
{\tiny\textcolor{Muted}{引用依据：本页综合三篇主讲材料的研究定位、证据等级和公开来源状态。}}
\end{frame}


\begin{frame}[t,shrink=6]{Arm Confidential Computing Architecture：内容摘要}
\DeckKicker{内容摘要}
\SlideMeta{证据等级}{E1 / 基础证据}{来源状态：本地 PDF 已核验}
{\large\bfseries\color{Ink} 提出并验证 Arm CCA 的核心设计，包括 Realm、RME、RMM、GPT/GPC 与 lifecycle\par}\vspace{1.8mm}
\begin{columns}[T,totalwidth=\textwidth]
  \begin{column}{0.45\textwidth}
    \fcolorbox{Rule}{Paper}{%
  \begin{minipage}[t]{0.97\linewidth}
    \vspace{1.1mm}
    \hspace{1.4mm}\begin{minipage}{0.92\linewidth}
      \PanelTitle{内容说明}
      \scriptsize \begin{itemize}
  \item 提出并验证 Arm CCA 的核心设计，包括 Realm、RME、RMM、GPT/GPC 与 lifecycle
  \item OSDI 2022 paper is the foundational peer-reviewed system treatment of Arm CCA, RME, and RMM.
\end{itemize}
    \end{minipage}
    \vspace{1mm}
  \end{minipage}}
  \end{column}
  \begin{column}{0.49\textwidth}
    \fcolorbox{Rule}{Panel}{%
  \begin{minipage}[t]{0.97\linewidth}
    \vspace{1.1mm}
    \hspace{1.4mm}\begin{minipage}{0.92\linewidth}
      \PanelTitle{证据定位}
      \scriptsize {\tiny
\begin{tabularx}{\linewidth}{@{}YY@{}}
\toprule
\textbf{字段} & \textbf{内容} \\
\midrule
论文定位 & 开创/基础入口 \\
材料类型 & 系统论文 \\
证据等级 & E1 / 基础证据 \\
来源状态 & 本地 PDF 已核验 \\
\bottomrule
\end{tabularx}
}
    \end{minipage}
    \vspace{1mm}
  \end{minipage}}
  \end{column}
\end{columns}
\vspace{1mm}
{\tiny\textcolor{Muted}{引用依据：引用依据为论文原文、官方规范或公开材料；本页结论按 E1 / 基础证据 的证据等级使用，不外推到未验证平台。}}
\end{frame}


\begin{frame}[t,shrink=6]{Arm Confidential Computing Architecture：研究背景}
\DeckKicker{研究背景}
\SlideMeta{证据等级}{E1 / 基础证据}{来源状态：本地 PDF 已核验}
{\large\bfseries\color{Ink} 云场景下 host 仍需管理资源，但不应看到 tenant workload 数据\par}\vspace{1.8mm}
\begin{columns}[T,totalwidth=\textwidth]
  \begin{column}{0.45\textwidth}
    \fcolorbox{Rule}{Paper}{%
  \begin{minipage}[t]{0.97\linewidth}
    \vspace{1.1mm}
    \hspace{1.4mm}\begin{minipage}{0.92\linewidth}
      \PanelTitle{内容说明}
      \scriptsize \begin{itemize}
  \item 云场景下 host 仍需管理资源，但不应看到 tenant workload 数据
\end{itemize}
    \end{minipage}
    \vspace{1mm}
  \end{minipage}}
  \end{column}
  \begin{column}{0.49\textwidth}
    \fcolorbox{Rule}{Panel}{%
  \begin{minipage}[t]{0.97\linewidth}
    \vspace{1.1mm}
    \hspace{1.4mm}\begin{minipage}{0.92\linewidth}
      \PanelTitle{问题边界}
      \scriptsize {\tiny
\begin{tabularx}{\linewidth}{@{}YY@{}}
\toprule
\textbf{维度} & \textbf{说明} \\
\midrule
保护对象 & Arm CCA / RME / RMM 基础架构 \\
传统瓶颈 & 围绕 Realm、RME、RMM、granule ownership 与 lifecycle 建立 Arm CCA 主线。 \\
本文入口 & Arm Confidential Computing Architecture \\
证据限制 & E1 / 基础证据 \\
\bottomrule
\end{tabularx}
}
    \end{minipage}
    \vspace{1mm}
  \end{minipage}}
  \end{column}
\end{columns}
\vspace{1mm}
{\tiny\textcolor{Muted}{引用依据：引用依据为论文原文、官方规范或公开材料；本页结论按 E1 / 基础证据 的证据等级使用，不外推到未验证平台。}}
\end{frame}


\begin{frame}[t,shrink=6]{Arm Confidential Computing Architecture：解决方案}
\DeckKicker{解决方案}
\SlideMeta{证据等级}{E1 / 基础证据}{来源状态：本地 PDF 已核验}
{\large\bfseries\color{Ink} 把资源管理和数据可见性分离，用硬件 granule ownership 和 RMM lifecycle 管理 Realm\par}\vspace{1.8mm}
\begin{columns}[T,totalwidth=\textwidth]
  \begin{column}{0.45\textwidth}
    \fcolorbox{Rule}{Paper}{%
  \begin{minipage}[t]{0.97\linewidth}
    \vspace{1.1mm}
    \hspace{1.4mm}\begin{minipage}{0.92\linewidth}
      \PanelTitle{内容说明}
      \scriptsize \begin{itemize}
  \item 把资源管理和数据可见性分离，用硬件 granule ownership 和 RMM lifecycle 管理 Realm
\end{itemize}
    \end{minipage}
    \vspace{1mm}
  \end{minipage}}
  \end{column}
  \begin{column}{0.49\textwidth}
    \fcolorbox{Rule}{Panel}{%
  \begin{minipage}[t]{0.97\linewidth}
    \vspace{1.1mm}
    \hspace{1.4mm}\begin{minipage}{0.92\linewidth}
      \PanelTitle{核心设计拆解}
      \scriptsize \begin{itemize}
  \item 01. Realm world and host/tenant separation
  \item 02. Granule ownership and GPT/GPC
  \item 03. RMM-mediated lifecycle and attestation path
\end{itemize}
    \end{minipage}
    \vspace{1mm}
  \end{minipage}}
  \end{column}
\end{columns}
\vspace{1mm}
{\tiny\textcolor{Muted}{引用依据：引用依据为论文原文、官方规范或公开材料；本页结论按 E1 / 基础证据 的证据等级使用，不外推到未验证平台。}}
\end{frame}


\begin{frame}[t,shrink=6]{Arm Confidential Computing Architecture：实验结果}
\DeckKicker{实验结果}
\SlideMeta{证据等级}{E1 / 基础证据}{来源状态：本地 PDF 已核验}
{\large\bfseries\color{Ink} OSDI 系统论文提供设计验证和原型评估\par}\vspace{1.8mm}
\begin{columns}[T,totalwidth=\textwidth]
  \begin{column}{0.45\textwidth}
    \fcolorbox{Rule}{Paper}{%
  \begin{minipage}[t]{0.97\linewidth}
    \vspace{1.1mm}
    \hspace{1.4mm}\begin{minipage}{0.92\linewidth}
      \PanelTitle{内容说明}
      \scriptsize \begin{itemize}
  \item OSDI 系统论文提供设计验证和原型评估
  \item 本阶段只记录可核对的证据类型，定量细节在 页面重写 中回引本地 PDF
\end{itemize}
    \end{minipage}
    \vspace{1mm}
  \end{minipage}}
  \end{column}
  \begin{column}{0.49\textwidth}
    \fcolorbox{Rule}{Panel}{%
  \begin{minipage}[t]{0.97\linewidth}
    \vspace{1.1mm}
    \hspace{1.4mm}\begin{minipage}{0.92\linewidth}
      \PanelTitle{实验/证据边界}
      \scriptsize {\tiny
\begin{tabularx}{\linewidth}{@{}YYY@{}}
\toprule
\textbf{对象} & \textbf{可支撑} & \textbf{边界} \\
\midrule
数据/来源 & Arm Confidential Computing Architecture & 本地 PDF 已核验 \\
Baseline/对照 & 以论文或规范原文为准 & 不跨材料外推 \\
指标/对象 & 只使用当前来源可证明的信息 & E1 / 基础证据 \\
使用边界 & 可进入本方向演进图 & 不能替代更强证据 \\
\bottomrule
\end{tabularx}
}
    \end{minipage}
    \vspace{1mm}
  \end{minipage}}
  \end{column}
\end{columns}
\vspace{1mm}
{\tiny\textcolor{Muted}{引用依据：引用依据为论文原文、官方规范或公开材料；本页结论按 E1 / 基础证据 的证据等级使用，不外推到未验证平台。}}
\end{frame}


\begin{frame}[t,shrink=6]{Arm Confidential Computing Architecture：文章评价}
\DeckKicker{文章评价}
\SlideMeta{证据等级}{E1 / 基础证据}{来源状态：本地 PDF 已核验}
{\large\bfseries\color{Ink} Arm CCA 主线入口，能把 Realm/RME/RMM 的安全边界讲清\par}\vspace{1.8mm}
\begin{columns}[T,totalwidth=\textwidth]
  \begin{column}{0.45\textwidth}
    \fcolorbox{Rule}{Paper}{%
  \begin{minipage}[t]{0.97\linewidth}
    \vspace{1.1mm}
    \hspace{1.4mm}\begin{minipage}{0.92\linewidth}
      \PanelTitle{内容说明}
      \scriptsize \begin{itemize}
  \item 设计优势：Arm CCA 主线入口，能把 Realm/RME/RMM 的安全边界讲清。
  \item 局限性：论文原型不等同于后续完整规范或量产实现。
  \item 商业化潜力：商业潜力强，因为机制进入 Arm 架构和 Linux/firmware 生态；风险在硬件普及、RMM 正确性和平台 attestation 集成。
  \item 证据边界：E1 / Foundational；本地 PDF 已核验。
\end{itemize}
    \end{minipage}
    \vspace{1mm}
  \end{minipage}}
  \end{column}
  \begin{column}{0.49\textwidth}
    \fcolorbox{Rule}{Panel}{%
  \begin{minipage}[t]{0.97\linewidth}
    \vspace{1.1mm}
    \hspace{1.4mm}\begin{minipage}{0.92\linewidth}
      \PanelTitle{文章评价}
      \scriptsize {\tiny
\begin{tabularx}{\linewidth}{@{}YY@{}}
\toprule
\textbf{维度} & \textbf{判断} \\
\midrule
设计优势 & Arm CCA 主线入口，能把 Realm/RME/RMM 的安全边界讲清。 \\
主要局限 & 论文原型不等同于后续完整规范或量产实现。 \\
商业落地 & 商业潜力强，因为机制进入 Arm 架构和 Linux/firmware 生态；风险在硬件普及、RMM 正确性和平台 attestation 集成。 \\
讲稿定位 & 开创/基础入口; 证据强度 E1 \\
\bottomrule
\end{tabularx}
}
    \end{minipage}
    \vspace{1mm}
  \end{minipage}}
  \end{column}
\end{columns}
\vspace{1mm}
{\tiny\textcolor{Muted}{引用依据：引用依据为论文原文、官方规范或公开材料；本页结论按 E1 / 基础证据 的证据等级使用，不外推到未验证平台。}}
\end{frame}


\begin{frame}[t,shrink=6]{Arm CCA 架构规范：内容摘要}
\DeckKicker{内容摘要}
\SlideMeta{证据等级}{E0 / Spec-standard SOTA}{来源状态：公开来源已核验，PDF 暂缺}
{\large\bfseries\color{Ink} 给出 CCA 架构接口、Realm、memory ownership 和 attestation 语义\par}\vspace{1.8mm}
\begin{columns}[T,totalwidth=\textwidth]
  \begin{column}{0.45\textwidth}
    \fcolorbox{Rule}{Paper}{%
  \begin{minipage}[t]{0.97\linewidth}
    \vspace{1.1mm}
    \hspace{1.4mm}\begin{minipage}{0.92\linewidth}
      \PanelTitle{内容说明}
      \scriptsize \begin{itemize}
  \item 给出 CCA 架构接口、Realm、memory ownership 和 attestation 语义
  \item First-party Arm CCA specification defines the current Realm model and standard evidence bas。
\end{itemize}
    \end{minipage}
    \vspace{1mm}
  \end{minipage}}
  \end{column}
  \begin{column}{0.49\textwidth}
    \fcolorbox{Rule}{Panel}{%
  \begin{minipage}[t]{0.97\linewidth}
    \vspace{1.1mm}
    \hspace{1.4mm}\begin{minipage}{0.92\linewidth}
      \PanelTitle{证据定位}
      \scriptsize {\tiny
\begin{tabularx}{\linewidth}{@{}YY@{}}
\toprule
\textbf{字段} & \textbf{内容} \\
\midrule
论文定位 & 代表性改进 \\
材料类型 & 规范/标准材料 \\
证据等级 & E0 / Spec-standard SOTA \\
来源状态 & 公开来源已核验，PDF 暂缺 \\
\bottomrule
\end{tabularx}
}
    \end{minipage}
    \vspace{1mm}
  \end{minipage}}
  \end{column}
\end{columns}
\vspace{1mm}
{\tiny\textcolor{Muted}{引用依据：引用依据为论文原文、官方规范或公开材料；本页结论按 E0 / Spec-standard SOTA 的证据等级使用，不外推到未验证平台。}}
\end{frame}


\begin{frame}[t,shrink=6]{Arm CCA 架构规范：研究背景}
\DeckKicker{研究背景}
\SlideMeta{证据等级}{E0 / Spec-standard SOTA}{来源状态：公开来源已核验，PDF 暂缺}
{\large\bfseries\color{Ink} 研究论文需要落到稳定的硬件/固件 ABI 和平台行为\par}\vspace{1.8mm}
\begin{columns}[T,totalwidth=\textwidth]
  \begin{column}{0.45\textwidth}
    \fcolorbox{Rule}{Paper}{%
  \begin{minipage}[t]{0.97\linewidth}
    \vspace{1.1mm}
    \hspace{1.4mm}\begin{minipage}{0.92\linewidth}
      \PanelTitle{内容说明}
      \scriptsize \begin{itemize}
  \item 研究论文需要落到稳定的硬件/固件 ABI 和平台行为
\end{itemize}
    \end{minipage}
    \vspace{1mm}
  \end{minipage}}
  \end{column}
  \begin{column}{0.49\textwidth}
    \fcolorbox{Rule}{Panel}{%
  \begin{minipage}[t]{0.97\linewidth}
    \vspace{1.1mm}
    \hspace{1.4mm}\begin{minipage}{0.92\linewidth}
      \PanelTitle{问题边界}
      \scriptsize {\tiny
\begin{tabularx}{\linewidth}{@{}YY@{}}
\toprule
\textbf{维度} & \textbf{说明} \\
\midrule
保护对象 & Arm CCA / RME / RMM 基础架构 \\
传统瓶颈 & 围绕 Realm、RME、RMM、granule ownership 与 lifecycle 建立 Arm CCA 主线。 \\
本文入口 & Arm CCA 架构规范 \\
证据限制 & E0 / Spec-standard SOTA \\
\bottomrule
\end{tabularx}
}
    \end{minipage}
    \vspace{1mm}
  \end{minipage}}
  \end{column}
\end{columns}
\vspace{1mm}
{\tiny\textcolor{Muted}{引用依据：引用依据为论文原文、官方规范或公开材料；本页结论按 E0 / Spec-standard SOTA 的证据等级使用，不外推到未验证平台。}}
\end{frame}


\begin{frame}[t,shrink=6]{Arm CCA 架构规范：解决方案}
\DeckKicker{解决方案}
\SlideMeta{证据等级}{E0 / Spec-standard SOTA}{来源状态：公开来源已核验，PDF 暂缺}
{\large\bfseries\color{Ink} 规范化 CCA/RME 相关对象、状态、接口和安全语义\par}\vspace{1.8mm}
\begin{columns}[T,totalwidth=\textwidth]
  \begin{column}{0.45\textwidth}
    \fcolorbox{Rule}{Paper}{%
  \begin{minipage}[t]{0.97\linewidth}
    \vspace{1.1mm}
    \hspace{1.4mm}\begin{minipage}{0.92\linewidth}
      \PanelTitle{内容说明}
      \scriptsize \begin{itemize}
  \item 规范化 CCA/RME 相关对象、状态、接口和安全语义
\end{itemize}
    \end{minipage}
    \vspace{1mm}
  \end{minipage}}
  \end{column}
  \begin{column}{0.49\textwidth}
    \fcolorbox{Rule}{Panel}{%
  \begin{minipage}[t]{0.97\linewidth}
    \vspace{1.1mm}
    \hspace{1.4mm}\begin{minipage}{0.92\linewidth}
      \PanelTitle{核心设计拆解}
      \scriptsize \begin{itemize}
  \item 01. Realm security state and PAS ownership
  \item 02. Memory delegation, RIPAS and lifecycle vocabulary
  \item 03. Attestation and platform interface semantics
\end{itemize}
    \end{minipage}
    \vspace{1mm}
  \end{minipage}}
  \end{column}
\end{columns}
\vspace{1mm}
{\tiny\textcolor{Muted}{引用依据：引用依据为论文原文、官方规范或公开材料；本页结论按 E0 / Spec-standard SOTA 的证据等级使用，不外推到未验证平台。}}
\end{frame}


\begin{frame}[t,shrink=6]{Arm CCA 架构规范：实验结果}
\DeckKicker{实验结果}
\SlideMeta{证据等级}{E0 / Spec-standard SOTA}{来源状态：公开来源已核验，PDF 暂缺}
{\large\bfseries\color{Ink} 规范/Survey，无新实验；Arm CCA 架构规范 只支撑术语、taxonomy、接口语义或证据边界。\par}\vspace{1.8mm}
\begin{columns}[T,totalwidth=\textwidth]
  \begin{column}{0.45\textwidth}
    \fcolorbox{Rule}{Paper}{%
  \begin{minipage}[t]{0.97\linewidth}
    \vspace{1.1mm}
    \hspace{1.4mm}\begin{minipage}{0.92\linewidth}
      \PanelTitle{内容说明}
      \scriptsize \begin{itemize}
  \item 本页不展示独立系统实验，也不把二手归纳写成一手结果。
  \item 可支撑：First-party Arm CCA specification defines the current Realm model and standard eviden。
  \item 证据等级：E0 / Spec-standard SOTA；来源状态：公开来源已核验，PDF 暂缺。
  \item 涉及具体平台、协议状态或数值时，转到原始论文、官方规范或厂商材料核验。
\end{itemize}
    \end{minipage}
    \vspace{1mm}
  \end{minipage}}
  \end{column}
  \begin{column}{0.49\textwidth}
    \fcolorbox{Rule}{Panel}{%
  \begin{minipage}[t]{0.97\linewidth}
    \vspace{1.1mm}
    \hspace{1.4mm}\begin{minipage}{0.92\linewidth}
      \PanelTitle{实验/证据边界}
      \scriptsize {\tiny
\begin{tabularx}{\linewidth}{@{}YYY@{}}
\toprule
\textbf{对象} & \textbf{可支撑} & \textbf{边界} \\
\midrule
数据/来源 & Arm CCA 架构规范 & 公开来源已核验，PDF 暂缺 \\
Baseline/对照 & 以论文或规范原文为准 & 不跨材料外推 \\
指标/对象 & 只使用当前来源可证明的信息 & E0 / Spec-standard SOTA \\
使用边界 & 可进入本方向演进图 & 不能替代更强证据 \\
\bottomrule
\end{tabularx}
}
    \end{minipage}
    \vspace{1mm}
  \end{minipage}}
  \end{column}
\end{columns}
\vspace{1mm}
{\tiny\textcolor{Muted}{引用依据：引用依据为论文原文、官方规范或公开材料；本页结论按 E0 / Spec-standard SOTA 的证据等级使用，不外推到未验证平台。}}
\end{frame}


\begin{frame}[t,shrink=6]{Arm CCA 架构规范：文章评价}
\DeckKicker{文章评价}
\SlideMeta{证据等级}{E0 / Spec-standard SOTA}{来源状态：公开来源已核验，PDF 暂缺}
{\large\bfseries\color{Ink} 评价：Arm CCA 架构规范 适合做 taxonomy/规范入口；规范/Survey，无新实验，不能替代一手系统证据。\par}\vspace{1.8mm}
\begin{columns}[T,totalwidth=\textwidth]
  \begin{column}{0.45\textwidth}
    \fcolorbox{Rule}{Paper}{%
  \begin{minipage}[t]{0.97\linewidth}
    \vspace{1.1mm}
    \hspace{1.4mm}\begin{minipage}{0.92\linewidth}
      \PanelTitle{内容说明}
      \scriptsize \begin{itemize}
  \item 设计优势：写 CCA 机制和状态语义时最权威。
  \item 局限性：不是性能、可用性或具体实现正确性的证据。
  \item 商业化潜力：商业价值取决于 Armv9 CCA 平台、firmware/RMM、OS 和 verifier 生态同步成熟。
  \item 规范/Survey，无新实验；具体平台结论转到一手材料核验。
\end{itemize}
    \end{minipage}
    \vspace{1mm}
  \end{minipage}}
  \end{column}
  \begin{column}{0.49\textwidth}
    \fcolorbox{Rule}{Panel}{%
  \begin{minipage}[t]{0.97\linewidth}
    \vspace{1.1mm}
    \hspace{1.4mm}\begin{minipage}{0.92\linewidth}
      \PanelTitle{文章评价}
      \scriptsize {\tiny
\begin{tabularx}{\linewidth}{@{}YY@{}}
\toprule
\textbf{维度} & \textbf{判断} \\
\midrule
设计优势 & 写 CCA 机制和状态语义时最权威。 \\
主要局限 & 不是性能、可用性或具体实现正确性的证据。 \\
商业落地 & 商业价值取决于 Armv9 CCA 平台、firmware/RMM、OS 和 verifier 生态同步成熟。 \\
讲稿定位 & 代表性改进; 证据强度 E0 \\
\bottomrule
\end{tabularx}
}
    \end{minipage}
    \vspace{1mm}
  \end{minipage}}
  \end{column}
\end{columns}
\vspace{1mm}
{\tiny\textcolor{Muted}{引用依据：引用依据为论文原文、官方规范或公开材料；本页结论按 E0 / Spec-standard SOTA 的证据等级使用，不外推到未验证平台。}}
\end{frame}


\begin{frame}[t,shrink=6]{Realm Management Monitor 规范：内容摘要}
\DeckKicker{内容摘要}
\SlideMeta{证据等级}{E0 / Spec-standard SOTA}{来源状态：公开来源已核验，PDF 暂缺}
{\large\bfseries\color{Ink} RMM spec 定义 Realm 管理软件接口、RMI/RSI 角色和 lifecycle\par}\vspace{1.8mm}
\begin{columns}[T,totalwidth=\textwidth]
  \begin{column}{0.45\textwidth}
    \fcolorbox{Rule}{Paper}{%
  \begin{minipage}[t]{0.97\linewidth}
    \vspace{1.1mm}
    \hspace{1.4mm}\begin{minipage}{0.92\linewidth}
      \PanelTitle{内容说明}
      \scriptsize \begin{itemize}
  \item RMM spec 定义 Realm 管理软件接口、RMI/RSI 角色和 lifecycle
  \item First-party RMM specification defines Realm management and lifecycle semantics for CCA syst。
\end{itemize}
    \end{minipage}
    \vspace{1mm}
  \end{minipage}}
  \end{column}
  \begin{column}{0.49\textwidth}
    \fcolorbox{Rule}{Panel}{%
  \begin{minipage}[t]{0.97\linewidth}
    \vspace{1.1mm}
    \hspace{1.4mm}\begin{minipage}{0.92\linewidth}
      \PanelTitle{证据定位}
      \scriptsize {\tiny
\begin{tabularx}{\linewidth}{@{}YY@{}}
\toprule
\textbf{字段} & \textbf{内容} \\
\midrule
论文定位 & 当前 SOTA/补充边界 \\
材料类型 & 规范/标准材料 \\
证据等级 & E0 / Spec-standard SOTA \\
来源状态 & 公开来源已核验，PDF 暂缺 \\
\bottomrule
\end{tabularx}
}
    \end{minipage}
    \vspace{1mm}
  \end{minipage}}
  \end{column}
\end{columns}
\vspace{1mm}
{\tiny\textcolor{Muted}{引用依据：引用依据为论文原文、官方规范或公开材料；本页结论按 E0 / Spec-standard SOTA 的证据等级使用，不外推到未验证平台。}}
\end{frame}


\begin{frame}[t,shrink=6]{Realm Management Monitor 规范：研究背景}
\DeckKicker{研究背景}
\SlideMeta{证据等级}{E0 / Spec-standard SOTA}{来源状态：公开来源已核验，PDF 暂缺}
{\large\bfseries\color{Ink} RMM 位于 CCA TCB 核心，状态机或内存管理错误会破坏 Realm 安全\par}\vspace{1.8mm}
\begin{columns}[T,totalwidth=\textwidth]
  \begin{column}{0.45\textwidth}
    \fcolorbox{Rule}{Paper}{%
  \begin{minipage}[t]{0.97\linewidth}
    \vspace{1.1mm}
    \hspace{1.4mm}\begin{minipage}{0.92\linewidth}
      \PanelTitle{内容说明}
      \scriptsize \begin{itemize}
  \item RMM 位于 CCA TCB 核心，状态机或内存管理错误会破坏 Realm 安全
\end{itemize}
    \end{minipage}
    \vspace{1mm}
  \end{minipage}}
  \end{column}
  \begin{column}{0.49\textwidth}
    \fcolorbox{Rule}{Panel}{%
  \begin{minipage}[t]{0.97\linewidth}
    \vspace{1.1mm}
    \hspace{1.4mm}\begin{minipage}{0.92\linewidth}
      \PanelTitle{问题边界}
      \scriptsize {\tiny
\begin{tabularx}{\linewidth}{@{}YY@{}}
\toprule
\textbf{维度} & \textbf{说明} \\
\midrule
保护对象 & Arm CCA / RME / RMM 基础架构 \\
传统瓶颈 & 围绕 Realm、RME、RMM、granule ownership 与 lifecycle 建立 Arm CCA 主线。 \\
本文入口 & Realm Management Monitor 规范 \\
证据限制 & E0 / Spec-standard SOTA \\
\bottomrule
\end{tabularx}
}
    \end{minipage}
    \vspace{1mm}
  \end{minipage}}
  \end{column}
\end{columns}
\vspace{1mm}
{\tiny\textcolor{Muted}{引用依据：引用依据为论文原文、官方规范或公开材料；本页结论按 E0 / Spec-standard SOTA 的证据等级使用，不外推到未验证平台。}}
\end{frame}


\begin{frame}[t,shrink=6]{Realm Management Monitor 规范：解决方案}
\DeckKicker{解决方案}
\SlideMeta{证据等级}{E0 / Spec-standard SOTA}{来源状态：公开来源已核验，PDF 暂缺}
{\large\bfseries\color{Ink} 用规范接口约束 Realm creation、memory transition、measurement、run/exit 和 attestati。\par}\vspace{1.8mm}
\begin{columns}[T,totalwidth=\textwidth]
  \begin{column}{0.45\textwidth}
    \fcolorbox{Rule}{Paper}{%
  \begin{minipage}[t]{0.97\linewidth}
    \vspace{1.1mm}
    \hspace{1.4mm}\begin{minipage}{0.92\linewidth}
      \PanelTitle{内容说明}
      \scriptsize \begin{itemize}
  \item 用规范接口约束 Realm creation、memory transition、measurement、run/exit 和 attestation 相关状态
\end{itemize}
    \end{minipage}
    \vspace{1mm}
  \end{minipage}}
  \end{column}
  \begin{column}{0.49\textwidth}
    \fcolorbox{Rule}{Panel}{%
  \begin{minipage}[t]{0.97\linewidth}
    \vspace{1.1mm}
    \hspace{1.4mm}\begin{minipage}{0.92\linewidth}
      \PanelTitle{核心设计拆解}
      \scriptsize \begin{itemize}
  \item 01. RMI/RSI interface boundary
  \item 02. Realm and REC lifecycle
  \item 03. Granule state transition and measurement hooks
\end{itemize}
    \end{minipage}
    \vspace{1mm}
  \end{minipage}}
  \end{column}
\end{columns}
\vspace{1mm}
{\tiny\textcolor{Muted}{引用依据：引用依据为论文原文、官方规范或公开材料；本页结论按 E0 / Spec-standard SOTA 的证据等级使用，不外推到未验证平台。}}
\end{frame}


\begin{frame}[t,shrink=6]{Realm Management Monitor 规范：实验结果}
\DeckKicker{实验结果}
\SlideMeta{证据等级}{E0 / Spec-standard SOTA}{来源状态：公开来源已核验，PDF 暂缺}
{\large\bfseries\color{Ink} 规范/Survey，无新实验；Realm Management Monitor 规范 只支撑术语、taxonomy、接口语义或证据边界。\par}\vspace{1.8mm}
\begin{columns}[T,totalwidth=\textwidth]
  \begin{column}{0.45\textwidth}
    \fcolorbox{Rule}{Paper}{%
  \begin{minipage}[t]{0.97\linewidth}
    \vspace{1.1mm}
    \hspace{1.4mm}\begin{minipage}{0.92\linewidth}
      \PanelTitle{内容说明}
      \scriptsize \begin{itemize}
  \item 本页不展示独立系统实验，也不把二手归纳写成一手结果。
  \item 可支撑：First-party RMM specification defines Realm management and lifecycle semantics for CC。
  \item 证据等级：E0 / Spec-standard SOTA；来源状态：公开来源已核验，PDF 暂缺。
  \item 涉及具体平台、协议状态或数值时，转到原始论文、官方规范或厂商材料核验。
\end{itemize}
    \end{minipage}
    \vspace{1mm}
  \end{minipage}}
  \end{column}
  \begin{column}{0.49\textwidth}
    \fcolorbox{Rule}{Panel}{%
  \begin{minipage}[t]{0.97\linewidth}
    \vspace{1.1mm}
    \hspace{1.4mm}\begin{minipage}{0.92\linewidth}
      \PanelTitle{实验/证据边界}
      \scriptsize {\tiny
\begin{tabularx}{\linewidth}{@{}YYY@{}}
\toprule
\textbf{对象} & \textbf{可支撑} & \textbf{边界} \\
\midrule
数据/来源 & Realm Management Monitor 规范 & 公开来源已核验，PDF 暂缺 \\
Baseline/对照 & 以论文或规范原文为准 & 不跨材料外推 \\
指标/对象 & 只使用当前来源可证明的信息 & E0 / Spec-standard SOTA \\
使用边界 & 可进入本方向演进图 & 不能替代更强证据 \\
\bottomrule
\end{tabularx}
}
    \end{minipage}
    \vspace{1mm}
  \end{minipage}}
  \end{column}
\end{columns}
\vspace{1mm}
{\tiny\textcolor{Muted}{引用依据：引用依据为论文原文、官方规范或公开材料；本页结论按 E0 / Spec-standard SOTA 的证据等级使用，不外推到未验证平台。}}
\end{frame}


\begin{frame}[t,shrink=6]{Realm Management Monitor 规范：文章评价}
\DeckKicker{文章评价}
\SlideMeta{证据等级}{E0 / Spec-standard SOTA}{来源状态：公开来源已核验，PDF 暂缺}
{\large\bfseries\color{Ink} 评价：Realm Management Monitor 规范 适合做 taxonomy/规范入口；规范/Survey，无新实验，不能替代一手系统证据。\par}\vspace{1.8mm}
\begin{columns}[T,totalwidth=\textwidth]
  \begin{column}{0.45\textwidth}
    \fcolorbox{Rule}{Paper}{%
  \begin{minipage}[t]{0.97\linewidth}
    \vspace{1.1mm}
    \hspace{1.4mm}\begin{minipage}{0.92\linewidth}
      \PanelTitle{内容说明}
      \scriptsize \begin{itemize}
  \item 设计优势：能把 RMM 作为 CCA TCB 核心讲清。
  \item 局限性：不证明任一 RMM 实现正确，且本地 PDF 不可用时不能写不可见细节。
  \item 商业化潜力：落地依赖 RMM 实现验证、firmware 更新、硬件支持和 verifier 能否检查关键状态。
  \item 规范/Survey，无新实验；具体平台结论转到一手材料核验。
\end{itemize}
    \end{minipage}
    \vspace{1mm}
  \end{minipage}}
  \end{column}
  \begin{column}{0.49\textwidth}
    \fcolorbox{Rule}{Panel}{%
  \begin{minipage}[t]{0.97\linewidth}
    \vspace{1.1mm}
    \hspace{1.4mm}\begin{minipage}{0.92\linewidth}
      \PanelTitle{文章评价}
      \scriptsize {\tiny
\begin{tabularx}{\linewidth}{@{}YY@{}}
\toprule
\textbf{维度} & \textbf{判断} \\
\midrule
设计优势 & 能把 RMM 作为 CCA TCB 核心讲清。 \\
主要局限 & 不证明任一 RMM 实现正确，且本地 PDF 不可用时不能写不可见细节。 \\
商业落地 & 落地依赖 RMM 实现验证、firmware 更新、硬件支持和 verifier 能否检查关键状态。 \\
讲稿定位 & 当前 SOTA/补充边界; 证据强度 E0 \\
\bottomrule
\end{tabularx}
}
    \end{minipage}
    \vspace{1mm}
  \end{minipage}}
  \end{column}
\end{columns}
\vspace{1mm}
{\tiny\textcolor{Muted}{引用依据：引用依据为论文原文、官方规范或公开材料；本页结论按 E0 / Spec-standard SOTA 的证据等级使用，不外推到未验证平台。}}
\end{frame}


\begin{frame}[t,shrink=6]{Arm CCA / RME / RMM 基础架构：方向总结}
\DeckKicker{方向总结}
\SlideMeta{证据等级}{方向综述}{来源状态：公开材料综合}
{\large\bfseries\color{Ink} Arm CCA / RME / RMM 基础架构 的结论是：三篇主讲材料共同给出技术演进，但仍要用证据等级约束每个结论。\par}\vspace{1.8mm}
\begin{columns}[T,totalwidth=\textwidth]
  \begin{column}{0.45\textwidth}
    \fcolorbox{Rule}{Paper}{%
  \begin{minipage}[t]{0.97\linewidth}
    \vspace{1.1mm}
    \hspace{1.4mm}\begin{minipage}{0.92\linewidth}
      \PanelTitle{内容说明}
      \scriptsize \begin{itemize}
  \item Arm Confidential Computing Architecture：开创/基础入口，E1 / 基础证据。
  \item Arm CCA 架构规范：代表性改进，E0 / Spec-standard SOTA。
  \item Realm Management Monitor 规范：当前 SOTA/补充边界，E0 / Spec-standard SOTA。
  \item 本方向结论：先用三篇主讲材料建立演进关系，再用证据等级控制机制、实验和商业化结论。
\end{itemize}
    \end{minipage}
    \vspace{1mm}
  \end{minipage}}
  \end{column}
  \begin{column}{0.49\textwidth}
    \fcolorbox{Rule}{Panel}{%
  \begin{minipage}[t]{0.97\linewidth}
    \vspace{1.1mm}
    \hspace{1.4mm}\begin{minipage}{0.92\linewidth}
      \PanelTitle{本方向方法对比}
      \scriptsize {\tiny
\begin{tabularx}{\linewidth}{@{}YYY@{}}
\toprule
\textbf{论文} & \textbf{定位} & \textbf{选择理由} \\
\midrule
Arm Confidential Computing Architecture & 开创/基础入口 & OSDI 2022 paper is the foundational peer-reviewed system treatment of Arm。 \\
Arm CCA 架构规范 & 代表性改进 & First-party Arm CCA specification defines the current Realm model and sta。 \\
Realm Management Monitor 规范 & 当前 SOTA/补充边界 & First-party RMM specification defines Realm management and lifecycle sema。 \\
\bottomrule
\end{tabularx}
}
    \end{minipage}
    \vspace{1mm}
  \end{minipage}}
  \end{column}
\end{columns}
\vspace{1mm}
{\tiny\textcolor{Muted}{引用依据：总结页综合三篇主讲材料的选择理由、评价字段和证据等级。}}
\end{frame}
